\phantomsection
\color{unidarkgreen}\section*{\centering{\large{ПЕРЕЧЕНЬ СОКРАЩЕНИЙ}}}
\addcontentsline{toc}{section}{Перечень сокращений}
{\color{white}\fontsize{0.1pt}{0.1pt}\selectfont<begABBRS>}
\color{black}
\begin{longtable}{>{\raggedright\arraybackslash}m{2cm}>{\raggedright\arraybackslash}m{0.5cm}>{\raggedright\arraybackslash}m{20cm}}
\endfirsthead\endhead\endfoot\endlastfoot
GOOSE & -- & Generic Object-Oriented Substation Event; \\
IRF & -- & Internal Relay Fault; \\
PPS & -- & Pulse Per Second; \\
PTP & -- & Precision Time Protocol (протокол точного времени); \\
SNTP & -- & Simple Network Time Protocol (простой сетевой протокол времени); \\
АВР & -- & автоматическое включение резерва; \\
АМТЗ & -- & аварийная максимальная токовая защита; \\
АПВ & -- & автоматическое повторное включение; \\
АПТ & -- & автоматика пожаротушения; \\
АСУ~ТП & -- & автоматизированная система управления технологическими процессами; \\
АТ & -- & автотрансформатор; \\
АУ & -- & автоматическое ускорение; \\
АУВ & -- & автоматика управления выключателем; \\
АЦП & -- & аналого-цифровой преобразователь; \\
БИ & -- & блок испытательный; \\
БКУ & -- & блок команд управления; \\
БНН & -- & блокировка при неисправности цепей напряжения; \\
БНТ & -- & блокировка при бросках намагничивающего тока; \\
ВН & -- & высшее напряжение; \\
ВО & -- & вторичная обмотка; \\
ВЧ & -- & высокочастотный; \\
ВЧПП & -- & высокочастотный приемо-передатчик; \\
ГЗ & -- & газовая защита; \\
ДЗ & -- & дистанционная защита; \\
ДЗО & -- & дифференциальная защита ошиновки; \\
ДО & -- & дочерние общества; \\
ДТЗ & -- & дифференциальная токовая защита; \\
ДТЗт & -- & дифференциальная токовая защита с торможением; \\
ДТО & -- & дифференциальная токовая отсечка; \\
ДТм & -- & датчик температуры масла; \\
ДТо & -- & датчик температуры обмотки; \\
ДФЗ & -- & дифференциально-фазная защита; \\
ЕЭС & -- & единая энергетическая система; \\
ЗДЗ & -- & защита от дуговых замыканий; \\
ЗИП & -- & запасные части, инструменты и принадлежности; \\
ЗНР & -- & защита от неполнофазного режима; \\
ЗНФ & -- & защита от непереключения фаз; \\
ЗП & -- & защита от перегрузки; \\
ЗПО & -- & защита от потери охлаждения; \\
ИО & -- & измерительный орган / избирательный орган; \\
ИЧМ & -- & интерфейс человек-машина; \\
КЗ & -- & короткое замыкание; \\
КИ & -- & контроль изоляции; \\
КОН & -- & контроль отсутствия напряжения; \\
КПОВ & -- & контроль перевода на обходной выключатель; \\
КПОН & -- & комбинированный пусковой орган напряжения; \\
КС & -- & контроль синхронизма / контрольная сумма; \\
КСВ & -- & контроль силового выключателя; \\
КСН & -- & контроль синхронизма и напряжения; \\
КЦН & -- & контроль цепей напряжения; \\
КЦТ & -- & контроль цепей тока; \\
ЛВ & -- & логика включения; \\
ЛД & -- & логика деления; \\
ЛО & -- & логика отключения; \\
ЛППТ & -- & логика пуска пожаротушения; \\
ЛРТ & -- & линейный регулировочный трансформатор; \\
ЛЭП & -- & линия электропередачи; \\
МТЗ & -- & максимальная токовая защита; \\
МЭК & -- & международная электротехническая комиссия; \\
НВЧЗ & -- & направленная высокочастотная защита; \\
НН & -- & низшее напряжение; \\
ОВ & -- & оперативный вывод / обходной выключатель; \\
ОЗУ & -- & оперативное запоминающее устройство; \\
ОК & -- & отсечной клапан; \\
ОНМ & -- & орган направления мощности; \\
ОТ & -- & оперативный ток; \\
ОУ & -- & оперативное ускорение; \\
ПА & -- & противоаварийная автоматика; \\
ПЗУ & -- & постоянное запоминающее устройство; \\
ПК & -- & предохранительный клапан; \\
ПО & -- & пусковой орган / программное обеспечение; \\
ПРД & -- & передатчик; \\
ПРМ & -- & приемник; \\
РД & -- & реле давления; \\
РЗ & -- & релейная защита; \\
РЗА & -- & релейная защита и автоматика; \\
РНМ & -- & реле направления мощности; \\
РПВ & -- & реле положения «включено»; \\
РПН & -- & устройство регулирования напряжения под нагрузкой; \\
РС & -- & регистрация событий / реле сопротивления; \\
РТ & -- & реле тока; \\
РТПО & -- & реле (орган) тока пуска охлаждения; \\
РУ & -- & распределительное устройство; \\
РФ & -- & реле фиксации; \\
РФК & -- & реле фиксации команд; \\
РФП & -- & реле фиксации включенного положения; \\
РЭ & -- & руководство по эксплуатации; \\
СВ & -- & секционный выключатель; \\
СН & -- & среднее напряжение; \\
СО & -- & система охлаждения; \\
СШ & -- & система (секция) шин; \\
ТЗ & -- & технологические защиты; \\
ТЗНП & -- & токовая защита нулевой последовательности; \\
ТЗО & -- & токовая защита ошиновки; \\
ТК & -- & токовый контроль / текущий контроль; \\
ТН & -- & трансформатор напряжения; \\
ТНЗНП & -- & токовая направленная защита нулевой последовательности; \\
ТО & -- & токовая отсечка; \\
ТС & -- & технологическая сигнализация; \\
ТТ & -- & трансформатор тока; \\
ТУ & -- & телеуправление / телеускорение / технические условия; \\
УВ & -- & управление выключателем / управляющее воздействие; \\
УРОВ & -- & устройство резервирования при отказе выключателя; \\
УС & -- & улавливание синхронизма; \\
ФПВ & -- & фиксация положения выключателя; \\
ФСУ & -- & функциональная схема устройства; \\
ЦОС & -- & цифровая обработка сигналов; \\
ЦП & -- & центральный процессор; \\
ЦУ & -- & цепи управления; \\
ШСВ & -- & шиносоединительный выключатель; \\
ШЭТ & -- & шкаф электротехнический типовой; \\
ЭМВ & -- & электромагнит включения; \\
ЭМО & -- & электромагнит отключения; \\
ЭМУ & -- & электромагнит управления. \\
{\color{white}\fontsize{0.1pt}{0.1pt}\selectfont<endABBRS>}
\end{longtable}
